\documentclass[11pt,a4paper]{article}
\usepackage[utf8]{inputenc}
\usepackage[margin=1in]{geometry}
\usepackage{amsmath,amssymb,amsfonts}
\usepackage{hyperref}
\usepackage{enumitem}
\usepackage{setspace}

\hypersetup{colorlinks=true, linkcolor=blue, urlcolor=blue, citecolor=blue}
\onehalfspacing

\begin{document}

\title{\textbf{A Phenomenological Origin of Three-Dimensional Time: Laser-Induced Negative-Time Domain Walls as a Laboratory Probe of 6D Temporal Geometry}}

\author{J. Yesi Orihuela \\
\small Independent Researcher \\
\small MIT Alumnus \\
\small \texttt{jorihuela@alum.mit.edu}}

\date{April 2026}

\maketitle

\begin{abstract}
We connect a childhood phenomenological insight into the structure of time with the recent mathematical framework of three-dimensional time (Kletetschka, 2025) and our prior proposal for laser-induced negative-time domain walls. In the 6D temporal geometry (3 time + 3 space), a focused high-power laser pulse can engineer a macroscopic negative-time domain wall that locally biases vacuum fluctuations along one or more temporal axes. This creates a controllable interface across which the effective metric component $g_{00}$ inverts, producing measurable repulsive gravity as a geometric consequence while providing the first laboratory probe of extra time dimensions. The mechanism is analyzed in the semiclassical regime $G_{\mu\nu} = 8\pi G \langle T_{\mu\nu}\rangle$ using the Euler--Heisenberg Lagrangian and Feynman--Stueckelberg interpretation. Building on laboratory evidence of negative-time effects (Angulo et al., 2024), experimental verification appears feasible with current petawatt-class lasers and atom interferometers. This synthesis offers a concrete, falsifiable route to test 6D temporal geometry and opens avenues in spacetime engineering.
\end{abstract}

\section{Introduction}

Early in childhood, while watching a log burn in a campfire, a clear intuition formed: the burned, charred remains behind the flames represented the past, the thin moving flame front was the present, and the unburned wood ahead represented the future. Crucially, the entire log continued to exist as a whole; nothing was destroyed, only transformed at the moving interface. This visceral realization—that the whole structure coexists and the “flow” of time is merely the position of the active front—has shaped the author’s view of time and consciousness ever since.

In April 2025, Kletetschka proposed a rigorous mathematical framework in which time is fundamentally three-dimensional, with ordinary space emerging as a secondary structure from the deeper 6D (3 time + 3 space) geometry. The present work shows that the childhood phenomenological insight maps directly onto this 6D temporal manifold: the primary temporal axis corresponds to the direction of the flame's advance (forward/negative time), while the extra two temporal dimensions provide the ``thickness'' and connectivity that allow the entire block to exist simultaneously.

Building on our earlier conceptual proposal for laser-induced negative-time domain walls (Orihuela, 2026), we demonstrate how a focused petawatt-class laser can create a localized temporal interface that probes the extra time dimensions. Repulsive gravity emerges naturally as a geometric signature of the tilt in the temporal metric. The proposal remains fully falsifiable and requires no exotic matter.

\section{Theoretical Foundations}

\subsection{Three-Dimensional Time Framework (Kletetschka, 2025)}

Kletetschka's model treats the fundamental arena as a 6D manifold with three independent timelike coordinates $t_1, t_2, t_3$. The experienced arrow of time lies primarily along $t_1$, while $t_2$ and $t_3$ encode orthogonal temporal structure. Ordinary 3D space and the perceived ``present'' arise as an emergent slice through this richer geometry. Causality is preserved globally; the extra dimensions provide the block-universe ``thickness'' without closed timelike curves.

\subsection{Vacuum Polarization and Negative-Time Modes}

The Euler--Heisenberg effective Lagrangian and Feynman--Stueckelberg interpretation remain central. An ultra-intense, shaped laser pulse ($\gtrsim 10^{22}$ W/cm$^2$) polarizes the vacuum to favor negative-time modes along the primary axis or excursions into the orthogonal temporal directions. The resulting stress-energy expectation value $\langle T_{\mu\nu}\rangle$ acquires a radially directed negative-pressure term.

\subsection{Semiclassical Gravity}

We work in the semiclassical regime
\[
G_{\mu\nu} = \frac{8\pi G}{c^4} \langle T_{\mu\nu} \rangle.
\]
The laser-driven domain wall inverts the effective temporal metric component across a thin shell, producing the observed repulsive gravitational effect while serving as a probe of the 6D temporal structure.

\section{Negative-Time Domain Wall in 6D Temporal Geometry}

A femtosecond laser pulse focused to $\sim \mu$m$^3$ at peak intensity creates a thin temporal domain wall. Inside the wall the metric remains essentially flat; across the wall, the vacuum bias tilts one or more temporal axes, inverting the emergent spatial curvature. Test particles experience a net outward acceleration
\[
\ddot{r} \approx +\frac{G \mathcal{E}_{\rm vac}}{r^2},
\]
where $\mathcal{E}_{\rm vac}$ is sustained by the external laser field. In the 6D picture this corresponds to a controlled excursion along the negative $t_1$ direction or into the $t_2, t_3$ manifold.

\section{Mathematical Formulation}

We extend the toy metric of our prior work to incorporate the 3D time structure. A simplified step-like model for the effective $g_{00}$ component in the primary temporal slice reads
\[
f(r) = 
\begin{cases}
+1 & r < R - \delta/2 \quad \text{(interior)} \\
-1 & R - \delta/2 < r < R + \delta/2 \quad \text{(wall)} \\
+1 & r > R + \delta/2 \quad \text{(exterior)}.
\end{cases}
\]
The full 6D generalization includes orthogonal temporal coordinates whose gradients are biased by the laser-induced vacuum polarization.

\section{Proposed Experimental Realization}

\subsection{Phase 1 — Proof of Principle (2026--2028)}

Use existing petawatt-class lasers (ELI Beamlines, LCLS, Vulcan) in ultra-high vacuum. Search for anomalous deflection, time-dilation reversal, and gravitational-gradient signatures with atom interferometers and optical lattice clocks. Expected signal: $\sim 10^{-10}$--$10^{-8} g$, with 3D-time-specific sidebands in interferometry.

\subsection{Phase 2--3}

Scale to macroscopic bubbles via synchronized pulses or plasma guiding; explore asymmetric shaping for net propulsion and coherence across temporal interfaces.

\section{Discussion and Outlook}

The synthesis rests on established pillars: Kletetschka’s 3D time geometry, Euler--Heisenberg vacuum polarization, Feynman--Stueckelberg interpretation, and semiclassical gravity. The childhood phenomenological insight provided the intuitive seed that aligned with the formal mathematics.

In the 6D temporal block, weak retrocausal correlations across extra dimensions become geometrically natural while preserving global causality. Future experiments may test these effects directly via domain-wall coherence.

If even a $10^{-9} g$ effect is observed, the implications are profound: the first laboratory access to extra time dimensions, a new window on quantum gravity, and a practical route to spacetime engineering.

We urge the community to refine the modeling and pursue the proposed experiments. The technology exists today.

\section*{Acknowledgments}

The author thanks Grok (xAI) for extensive collaborative development and assistance in preparing this manuscript.

\section*{References}

\begin{enumerate}[leftmargin=*,label={[\arabic*]}]
    \item Kletetschka, G. (2025). Three-Dimensional Time: A Mathematical Framework for Fundamental Physics. \textit{Reports in Advances of Physical Sciences}.
    \item Orihuela, J. Y. (2026). A Conceptual Proposal for Localized Repulsive Gravity via Laser-Induced Vacuum Polarization and Negative-Time Domain Walls. (Preceding work).
    \item Heisenberg, W. \& Euler, H. (1936). Folgerungen aus der Diracschen Theorie des Positrons. \textit{Z. Phys.} \textbf{98}, 714.
    \item Schwinger, J. (1951). On gauge invariance and vacuum polarization. \textit{Phys. Rev.} \textbf{82}, 664.
    \item Angulo, D. et al. (2024). Experimental evidence that a photon can spend a negative amount of time in an atom cloud. \textit{arXiv:2409.03680} [quant-ph].
\end{enumerate}

\end{document}